Let @@M0@@ be a locally free @@M1@@-Module of finite rank @@M2@@. There exists a torsor @@M3@@ under @@M4@@ such that @@M5@@, where @@M6@@ is regarded as a @@M7@@-Module by the standard representation. We can then define a @@M8@@-homomorphism from @@M9@@ to @@M10@@ by associating to every representation @@M11@@ of @@M12@@ in a free @@M13@@-module of rank @@M14@@ the class of the locally free @@M15@@-Module
@@M16@@, where @@M17@@ is regarded as a @@M18@@-module by the representation @@M19@@. Under this homomorphism, the class of @@M20@@ is therefore the image of the standard representation. To check a polynomial relation among the @@M21@@, it therefore suffices that this relation hold for the standard representation in @@M22@@.

We can argue similarly with two locally free @@M23@@-Modules @@M24@@ of ranks @@M25@@; we can then find a torsor @@M26@@ under @@M27@@ such that @@M28@@, where @@M29@@ is regarded as a @@M30@@-Module by letting @@M31@@ act by the standard representation and @@M32@@ by the trivial representation, and conversely for @@M33@@. We then define a @@M34@@-homomorphism from @@M35@@ to @@M36@@ in the same way as above, and are thus reduced to checking the relations among the @@M37@@, @@M38@@ for the corresponding elements of @@M39@@.

We are therefore reduced to showing that, for every group scheme @@M40@@ over @@M41@@ which is a finite product of linear group schemes, the pre-@@M42@@-ring @@M43@@ is a @@M44@@-ring. For this, we can refer to [5], where @@M45@@ is computed explicitly. To avoid using this result, we can reason in a manner analogous to 3.2, using 1.14.

